% [原文 p.407]
\chapter{WKB cameral 网络}

\ChapterAbstract{本章我们关心的是把非阿贝尔化映射的研究推广到一般复约化李群 $G$ 的 $G$-局部系统。进行这一重要推广所需的第一件工具是 cameral 覆盖，它在本项目的语境下有效地推广了谱覆盖的概念。由此引出一种新的网络构造，称为 WKB cameral 网络。}

Gaiotto、Moore 与 Neitzke [154] 的非阿贝尔化映射（我们在第 15 章考察过）已由 Ionita [214] 与 Morrissey [274] 的学位论文从 $\SL_n(\C)$ 与 $\GL_n(\C)$ 的情形推广到任意复约化李群 $G$（另见他们的合著 [215]）。cameral 覆盖这一概念由 Donagi [112] 首次引入，它给出谱覆盖的推广，进而给出任意约化群情形下 Hitchin 系统的推广（[124] 与 [311]）。一般群 $G$ 的非阿贝尔化映射是沿一个基本谱网络经由“切割与重粘”过程产生的代数映射，从而确定一个代数叠之间的态射：源是某类带额外结构的局部系统在一个开子集上的模空间，该开子集位于非紧（带穿孔）黎曼曲面 $X$ 的 cameral 覆盖 $\tilde X\to X$ 之中；靶是 $X$ 上 $G$-局部系统的模空间。该映射在 $G=\GL_n(\C)$ 与 $G=\SL_n(\C)$ 时与 [154] 中的映射一致；但与 [154] 的构造不同，它不依赖于对局部系统或对分歧覆盖 $\tilde X\to X$ 的平凡化选择。此外，一般 $G$-局部系统的非阿贝尔化映射被猜想为局部同构；$G=\GL_n(\C)$ 情形的证据在 [154] 中于物理的严格性水平上给出。借助 [129]、[155]、[205] 的工作，该命题对 $G=\SL_2(\C)$ 在粗模空间的层面上已知成立。只要该猜想成立，该映射就给出 $G$-局部系统空间上的一个 étale 坐标卡。本章这个短章的目的是引入 cameral 覆盖 $\tilde X\to X$ 的概念，并解释 WKB cameral 网络这一工具：它由 $\tilde X$ 上一些定向曲线段组成，每条曲线段都用 $\g$ 的一个根作标记。此外，我们描述这些 Stokes 曲线如何下降为 $X$ 上的一组定向曲线，从而得到 $X$ 上的一个谱网络。关于一般 $G$-局部系统的非阿贝尔化映射的构造，
% [原文 p.408]
以及支持“该映射是局部同构”这一猜想的证据，我们请读者参阅 Ionita 与 Morrissey 的原始工作 [214]、[274] 与 [215]。

\section{cameral 覆盖}

谱网络最初是为 A 型群引入的，并借助黎曼曲面 $X$ 上的谱覆盖 $\tilde X\to X$ 推广到 D 型与 E 型群。为了把理论推广到更一般的约化李群 $G$，首要步骤是引入 $X$ 上的一个分歧覆盖，它不依赖于表示 $\rho:G\to\GL_n(\C)$ 的选择。该覆盖带有 $G$ 的 Weyl 群在纤维上的作用，该作用在分歧点之外自由且传递。我们下面把这一点说得更精确。

设 $G$ 是复约化李群，相应的复李代数为 $\g$；设 $\mathcal{E}$ 是闭、连通、可定向、亏格 $g\ge2$ 的黎曼曲面 $X$ 上的主 $G$-丛。又记 $\operatorname{ad}(\mathcal{E}):=\mathcal{E}\times_G\g$ 为 $\mathcal{E}$ 的伴随丛，它是 $X$ 上的向量丛。这里的扭乘积 $\mathcal{E}\times_G\g$ 是乘积丛 $\mathcal{E}\times\g$ 按等价关系 $(e\cdot g,x)\sim(e,\operatorname{ad}_g(x))$ 所得的商，其中 $e$ 取遍 $\mathcal{E}$ 的截面，$g\in G$、$x\in\g$ 任意。

\begin{definition}[25.1] 对如上 $X$ 与 $G$，固定 $X$ 上的一个线丛 $L$。$X$ 上 $L$-扭的 $G$-Higgs 丛是一个二元组 $(\mathcal{E},\varphi)$，其中 $\mathcal{E}$ 是 $X$ 上的主 $G$-丛，而 $\varphi$ 是 $\Gamma(X,\operatorname{ad}(\mathcal{E})\otimes L)$ 中的一个截面；截面 $\varphi$ 称为 Higgs 场。\end{definition}

\begin{remark}[25.2] 对群表示 $\rho:G\to\GL(V)$，可以考虑 $V$ 关于极大环面 $T\subset G$ 的权空间分解 $V=\bigoplus_{\lambda\in\Lambda}V_\lambda$，其中 $\Lambda$ 是 $\rho$ 的一族权（取自 $G$ 的权格）。现在可以仿照定义 8.11，考虑由 $\rho$ 的权的 $W$-轨道所指标的谱覆盖 $\tilde X_\lambda$，进而把所得的谱覆盖 $\tilde X_\rho\to X$ 视为这些子覆盖 $\tilde X_\lambda$ 的并。然而一个自然的问题是：这样的覆盖究竟在多大程度上依赖于表示 $\rho$？一般而言，存在无穷多个互不同构的这类子覆盖 $X_\lambda$，它们落在有限多个双有理等价类之中。Donagi 却构造了一个支配上述所有覆盖的覆盖 $\tilde X$。我们下面回顾这一构造，原始版本与完整讨论见 [112, Sect. 2]。\end{remark}

% [Tutor Note: 本页原文两处记号不一致：前文谱覆盖一律写作 $\tilde X_\lambda$、$\tilde X_\rho$（带波浪号），随后“there are infinitely many non-isomorphic such sub-covers $X_\lambda$”却漏掉了波浪号。按上下文（它们是 $\tilde X_\lambda$ 的同类对象）应作 $\tilde X_\lambda$；照原文译出并在此指出。]

固定 $G$ 的一个极大环面 $T\subset G$，并记 $N$ 为 $T$ 在 $G$ 中的正规化子。记 $W\cong N/T$ 为 Weyl 群，$\mathfrak{t}$ 为 $T$ 相应的李代数。Chevalley 定理断言，限制映射
\[
\C[\g]^G\to\C[\mathfrak{t}]^W
\]
是从李代数 $\g$ 上 $\operatorname{Ad}$-不变多项式函数到 Cartan 子代数 $\mathfrak{t}\subset\g$ 上 $W$-不变多项式函数的同构。随后我们考虑单射环同态
% [原文 p.409]
\[
\C[\mathfrak{t}]^W\cong\C[\g]^G\hookrightarrow\C[\g],
\]
并对这些环取素谱，从而定义仿射簇之间的一个满的 $G$-不变态射
\begin{equation}
\g\twoheadrightarrow\mathfrak{t}/W.
\tag{25.1}
\end{equation}
于是，与商映射 $\mathfrak{t}\to\mathfrak{t}/W$ 取纤维积，得到
\begin{equation}
\tilde{\g}:=\g\times_{\mathfrak{t}/W}\mathfrak{t},
\tag{25.2}
\end{equation}
而投影 $\pi:\tilde{\g}\to\g$ 是一个有限 $W$-Galois 态射。它称为李代数 $\g$ 的 cameral 覆盖。正则半单元素 $g\in\g$ 的纤维 $\pi^{-1}(g)$ 与 $\mathfrak{t}^*$ 中的房（chamber）之集等同。

现在设 $K_X=T^*X$ 为 $X$ 上的典范线丛，$D$ 为 $X$ 中有限多个点构成的既约有效除子。我们也将把 $K_X(D)$ 记作线丛 $K_X\otimes\OO(D)$。上面的态射 (25.1) 是 $G$-不变且 $\C^*$-等变的，因此，对给定的 $K_X(D)$-扭 $G$-Higgs 丛 $(\mathcal{E},\varphi)$，该态射可延拓为态射
\[
|\operatorname{ad}(\mathcal{E})\otimes K_X(D)|\to|\mathfrak{t}\otimes K_X(D)|/W,
\]
这里 $|\cdot|$ 表示丛的全空间。现在可以如 (25.2) 那样与 $\mathfrak{t}\otimes K_X(D)$ 取纤维积，并经由 Higgs 场 $\varphi$ 拉回，从而得到下述定义：

\begin{definition}[25.3] 由 $K_X(D)$-扭 $G$-Higgs 丛 $(\mathcal{E},\varphi)$ 确定的 cameral 覆盖由映射 $\pi:\tilde X\to X$ 给出，其中
\[
\tilde X:=\varphi^*\left(|\operatorname{ad}(\mathcal{E})\otimes K_X(D)|\times_{|\mathfrak{t}\otimes K_X(D)|/W}|\mathfrak{t}\otimes K_X(D)|\right),
\]
而 $\pi$ 是到第一个因子的投影。\end{definition}

更具体地说，截面空间
\[
\mathcal{B}(X,G,K_X(D)):=\Gamma(X,\mathfrak{t}\otimes K_X(D)/W)
\]
中的一点 $b$（该空间称为 Hitchin 基）确定一个态射 $X\xrightarrow{b}\mathfrak{t}\otimes K_X(D)/W$。于是，联系于点 $b$ 的 cameral 覆盖 $\tilde X_b$ 是下图中的纤维积：
\[
\begin{array}{ccc}
\tilde X_b & \xrightarrow{\ \tilde b\ } & \mathfrak{t}\otimes K_X(D)\\[5pt]
{\scriptstyle \pi}\Big\downarrow & & \Big\downarrow\\[5pt]
X & \xrightarrow{\ b\ } & \mathfrak{t}\otimes K_X(D)/W .
\end{array}
\]

只要除子 $D$ 至少有 3 个互不相同的点，Hitchin 基 $\mathcal{B}(X,G,K_X(D))$ 就含有一个非空的 Zariski 开子集 $\mathcal{B}^{\diamond}(X,G,K_X(D))$，使得 $b\in\mathcal{B}^{\diamond}(X,G,K_X(D))$ 当且仅当 cameral 覆盖 $\tilde X_b$ 光滑。此外，
% [原文 p.410]
对这些点 $b\in\mathcal{B}^{\diamond}(X,G,K_X(D))$，$\tilde X_b$ 连通。cameral 覆盖 $\pi:\tilde X_b\to X$ 在分歧轨迹之外是 $X$ 上的主 $W$-丛。对所有 $b\in\mathcal{B}^{\diamond}(X,G,K_X(D))$，cameral 覆盖 $\pi:\tilde X_b\to X$ 的每个分歧点的阶都是 2；关于这些断言的更多细节见 [286, Sects. 4.6 and 4.7]。我们在此记 $P\subset X$ 为分支点集，$R\subset\tilde X_b$ 为映射 $\pi$ 的分歧点集。

\begin{example}[25.4] 当 $G=\GL_n(\C)$ 时，Weyl 群 $W$ 是对称群 $S_n$。此时 cameral 覆盖 $\tilde X_b\to X$ 是 $X$ 的 $n!$ 次覆盖，而此情形下的谱覆盖 $\tilde Y_b$ 次数为 $n$。差别在于：这里的 $\tilde Y_b$ 参数化特征多项式的特征值，而 $\tilde X_b$ 参数化这些特征值的排序。当 Hitchin 基中的点 $b$ 取得足够一般，并令 $S_{n-1}$ 为其中一个特征值的稳定化子时，上述两类覆盖（谱覆盖与 cameral 覆盖）由
\[
\tilde X_b/S_{n-1}\cong\tilde Y_b
\]
相联系。

注意当 $n=2$ 时，谱覆盖与 cameral 覆盖同构。\end{example}

\begin{remark}[25.5] 承接上例，我们指出：一般而言，cameral 覆盖在底曲面上依次支配根覆盖与谱覆盖（见 [112, Lem. 3.3] 以及第 13.6 节的图 13.19）。不过，对群 $\GL_n(\C)$ 在 $n=2$ 或 $n=3$ 的情形，根覆盖与 cameral 覆盖一致。\end{remark}

% [Tutor Note: 可用维数/次数核对上例：$\tilde{\g}\to\g$ 是 $|W|$ 次覆盖，$\GL_n(\C)$ 的 $|W|=n!$，与 $\tilde X_b\to X$ 的次数 $n!$ 一致；而在 $n=2$ 时 $|W|=2=n$，故谱覆盖与 cameral 覆盖次数相同且（由例中公式 $\tilde X_b/S_1\cong\tilde Y_b$）同构。另注意 $S_{n-1}\subset S_n$ 的阶为 $(n-1)!$，故 $\tilde X_b/S_{n-1}\to X$ 的次数为 $n!/(n-1)!=n$，恰与谱覆盖次数相符——这是例 25.4 中公式的一个快捷核对。]

\section{WKB cameral 网络}

在第 13.6 节我们看到了 WKB 谱网络的构造，它最初是在群 $G$ 为 $\SL_n(\C)$、$\GL_n(\C)$ 或 ADE 型时引入的（沿袭 [154] 与 [249]）。下面我们沿袭 Ionita 与 Morrissey 近期的工作 [214]、[274]、[215]，概述该构造对任意约化群 $G$ 的推广。

设 $G$ 是复约化李群，$\Delta$ 记相应李代数 $\g$ 的根集。固定 $\g$ 的一个 Chevalley 基后，每个根 $\alpha\in\Delta$ 都确定一个 $\mathfrak{sl}_2$-三元组。现在记 $s_\alpha\in W$ 为关于根超平面 $H_\alpha$ 的反射，并记 $p$ 为商映射 $N\to T$，其中 $N$ 是 $G$ 中极大环面 $T\subset G$ 的正规化子，Weyl 群 $W\cong N/T$。

Chevalley 基对每个根 $\alpha\in\Delta$ 确定一个元素 $n_\alpha\in N$，事实上 $n_\alpha\in p^{-1}(s_\alpha)$。现在我们可以引入下述概念：

\begin{definition}[25.6] 设 $b\in\mathcal{B}^{\diamond}(X,G,K_X(D))$。对任意点 $d\in D$，可以写成
\[
b(x)=\left(\sum_{i=1}^{\infty}a_ix^i\right)\mathrm{d}x,
\]
% [原文 p.411]
其中 $a_i\in\mathfrak{t}/W$，而 $x$ 是 $d$ 附近的局部坐标。我们把 $b$ 在 $d$ 处的留数定义为 $a_{-1}$ 的任意提升 $\tilde a_{-1}\in\mathfrak{t}$；它们不依赖于围绕点 $d$ 所选取的局部坐标 $x$。若对每个点 $d\in D$ 都有
\[
\tilde a_{-1}\in\mathfrak{t}\setminus\bigcup_{\alpha\in\Delta}H_\alpha,
\]
则称 $b$ 满足条件 $\mathbf{R}$。\end{definition}

注意上述条件 $\mathbf{R}$ 是 $W$-不变的，因此只要一个留数满足它，所有留数就都满足。此外，条件 $R$ 表明 cameral 覆盖 $\pi$ 在除子 $D$ 处非分歧。

记 $\mathcal{B}_R^{\diamond}(X,G,K_X(D))$ 为 $\mathcal{B}^{\diamond}(X,G,K_X(D))$ 中满足条件 $\mathbf{R}$ 的点 $b$ 所成的子集；这是一个 Zariski 开子集。现在可以给这样的点 $b$ 联系某些亚纯二次微分：

\begin{proposition}[25.7（{{{{{[214]}}}}, {{{{[274]}}}}, {{{{[215]}}}}}）] 设 $b\in\mathcal{B}_R^{\diamond}(X,G,K_X(D))$，并设 $\pi:\tilde X_b\to X$ 为相应的 $X$ 的光滑 cameral 覆盖。对每个根 $\alpha\in\Delta$，令 $\tilde X_{b,\alpha}:=\tilde X_b/\langle s_\alpha\rangle$，并考虑映射 $\pi$ 的如下分解
\[
\tilde X_b\xrightarrow{\pi_\alpha}\tilde X_{b,\alpha}\xrightarrow{p_\alpha}X.
\]
则 $b$ 确定 $\tilde X_{b,\alpha}$ 上的一个亚纯二次微分，它在 $\pi_\alpha$ 的分支点处有单零点，在 $D$ 的原像处有二阶极点。\end{proposition}

把由此结果得到的微分记为
\[
\omega_{b,\alpha}\in\Gamma\left(\tilde X_{b,\alpha},\left(K_{\tilde X_{b,\alpha}}(p_\alpha^*D)\right)^{\otimes 2}\right),
\]
并记
\[
\chi_{b,\alpha}\in\Gamma\left(\tilde X_b,\left(K_{\tilde X_{b,\alpha}}(p_\alpha^*D)\right)\right)
\]
为 $\tilde X_b$ 上的线性微分，它们由如下映射的复合得到
\[
\tilde X_b\xrightarrow{\tilde b}\mathfrak{t}_{\pi^*K_X(D)}\to\mathfrak{t}_{\pi_\alpha^*(K_{\tilde X_{b,\alpha}}(p_\alpha^*D))}\xrightarrow{\alpha}\pi_\alpha^*(K_{\tilde X_{b,\alpha}}(p_\alpha^*D)).
\]
这些线性微分各自确定 $\tilde X_b$ 上的一个定向叶状结构，其叶为满足
\[
\int_{t_0}^{t_1}\chi_{b,\alpha}(\gamma(s))\in\R^+
\]
的曲线 $\gamma:\R\to\tilde X_b$。我们称微分 $\chi_{b,\alpha}$ 的定向轨道为该叶状结构的极大叶。现在终于可以给出一般约化群 $G$ 设定下 Stokes 曲线的构造了：

\begin{definition}[25.8（{{{{{[214]}}}}, {{{{[274]}}}}, {{{{[215]}}}}}）] 对点 $b\in\mathcal{B}_R^{\diamond}(X,G,K_X(D))$，我们定义 WKB 构造 $\widetilde{\mathcal{W}}_b$，它由 cameral 覆盖 $\tilde X_b$ 上若干称为 Stokes 曲线的定向曲线段组成。这些 Stokes 曲线用 $\Delta$ 中的根作标记，并按如下步骤算法式地构造：
% [原文 p.412]
\begin{enumerate}
  \item 主 Stokes 曲线是 $\tilde X_b$ 上某个线性微分 $\chi_{b,\alpha}$ 的一条临界定向轨道。从映射 $\pi$ 的每个分歧点出发，恰有六条主 Stokes 曲线。
  \item 设 $x\in\tilde X_b$ 是若干 Stokes 曲线的一个孤立交点，这些曲线用某个凸根集 $C_{in}\subset\Delta$ 中互不相同的根作标记。则对限制凸包中的每个根 $\beta\in\operatorname{Conv}^{\N}_{C_{in}}$，有一条次 Stokes 曲线 $g_\beta$ 从交点 $x$ 出发。这条次 Stokes 曲线是由 $\chi_{b,\beta}$ 确定的定向叶状结构从 $x$ 出发的唯一叶。
  \item 每当已有两条或更多 Stokes 曲线相交时，就递归地产生新的次 Stokes 曲线。见图 25.1。
\end{enumerate}
\end{definition}

我们记 $J$ 为 Stokes 曲线的交点之集，称为 WKB 构造的关节。

% [Tutor Note: 定义 25.6 中的展开 $\sum_{i=1}^{\infty}a_ix^i$ 与随后使用的 $a_{-1}$ 自相矛盾：若求和从 $i=1$ 开始，则 $a_{-1}$ 不存在。由于 $b$ 取自 $\mathcal{B}^{\diamond}=\Gamma(X,\mathfrak{t}\otimes K_X(D)/W)$，它在 $D$ 处至多有一阶极点，故 Laurent 展开应为 $\sum_{i=-1}^{\infty}a_ix^i$，$a_{-1}$ 正是留数项。这是原文的一处笔误（照译并在此指出），核对时以 $b$ 的定义（$K_X(D)$-扭）为准。]

\Fig[0.72\linewidth]{figures/fig_25_1.png}{25.1}{Stokes 曲线（实线）是某些定向轨道（透明线）的线段。主 Stokes 曲线 $g_{\alpha_1}$ 与 $g_{\alpha_2}$ 在点 $x$ 相交，随后次 Stokes 曲线 $g_{\beta_1}$ 与 $g_{\beta_2}$ 从 $x$ 发出。}

WKB 构造 $\widetilde{\mathcal{W}}_b$ 可以看成一个有向图，其顶点集为 $R\cup J$，并对 $\widetilde{\mathcal{W}}_b\setminus\{R\cup J\}$ 的每个连通分支取一条有向边。在这个意义上我们作如下定义：

\begin{definition}[25.9] 若下列条件成立，则称 $\widetilde{\mathcal{W}}_b$ 是容许的：
\begin{itemize}
  \item cameral 覆盖 $\tilde X_b\setminus\pi^{-1}(D)$ 中没有任何点是关节集 $J$ 的聚点。
  \item Stokes 曲线不在 $\tilde X_b$ 的分歧点处相交，换言之 $R\cap J=\varnothing$。
% [原文 p.413]
  \item 对应于 $\widetilde{\mathcal{W}}_b$ 的有向图无环。
  \item 对每个关节，入射 Stokes 曲线之标记所成的根集 $C_{in}\subset\Delta$ 是凸的。
\end{itemize}
\end{definition}

下面给出 Ionita 与 Morrissey 意义下的 WKB cameral 网络的定义：

\begin{definition}[25.10（{{{{{[214]}}}}, {{{{[274]}}}}, {{{{[215]}}}}}）] 若存在子集 $X'\subset X\setminus D$，使得包含映射 $X'\hookrightarrow X\setminus D$ 是同伦等价，且 $\widetilde{\mathcal{W}}_b$ 在 $\tilde X_b':=\pi^{-1}(X')\subset\tilde X_b$ 上的限制只含有限多条 Stokes 曲线，则称容许的 WKB 构造 $\widetilde{\mathcal{W}}_b$ 为一个 WKB cameral 网络。\end{definition}

Ionita 与 Morrissey 的一个重要猜想断言：使 WKB 构造 $\widetilde{\mathcal{W}}_b$ 容许的点 $b\in\mathcal{B}_R^{\diamond}(X,G,K_X(D))$ 之集在经典拓扑中开且稠密。事实上，任何容许的 $\widetilde{\mathcal{W}}_b$ 都确实是 WKB cameral 网络，也就是说，容许性条件保证上面定义的 WKB 构造在 $\pi^{-1}(D)$ 的一个邻域之外总只有有限多条 Stokes 曲线。

最后，从 cameral 覆盖 $\pi:\tilde X_b\to X$ 上的 WKB cameral 网络 $\widetilde{\mathcal{W}}_b$ 可联系 $X$ 上的一个 WKB 谱网络 $\mathcal{W}_b$，它定义为 $X$ 上定向带标记曲线之并：$\mathcal{W}_b=\pi_*\widetilde{\mathcal{W}}_b$；Stokes 曲线 $g\in\mathcal{W}_b$ 的标记是 $\Hom_W(\tilde X_b|_g,\Delta)$ 的一个局部常值截面 $\psi_g$，它把曲线 $g$ 的任一原像 $\tilde g$ 映到给 $\tilde g$ 作标记的那个根。

对 WKB cameral 网络，从覆盖映射 $\pi:\tilde X_b\to X$ 的分歧点 $r\in R$ 出发有六条 Stokes 曲线。映射 $\pi$ 在点 $r$ 的一个小邻域上的限制次数为 2；又因 WKB cameral 网络是 $W$-等变的，相应的谱网络从每个分支点出发有三条 Stokes 曲线。这一现象与 [154] 和 [249] 中谱网络的构造一致。

% [Tutor Note: “六个分歧点上的 Stokes 曲线”与“分支点上的三条”并不矛盾：$\tilde X_b\to X$ 在分歧点 $r$ 附近是二次分歧覆盖（局部次数 2），而构造是 $W$-等变的，故 $\tilde X_b$ 上的六条主 Stokes 曲线两两被分歧覆盖认同，在 $X$ 上恰剩三条。核对方法：数 $\pi^{-1}(b_i)$ 与 $\pi_*$ 的重数——$\deg(\pi_*)=1$ 逐曲线计，局部按 $2:1$ 合并。]

% [原文 p.414]
% 本页（印刷 p.414）为章末空白页，原文无正文，故此处不置内容。

% [原文 p.415]
\chapter{共形 Virasoro 块的非阿贝尔化}

\ChapterAbstract{这里我们简要纳入 Hao 与 Neitzke [188] 关于谱网络的一项近期应用，即共形块的非阿贝尔化。Hao 与 Neitzke 在黎曼曲面 $C$ 上 $c=1$ 的 Virasoro 共形块与 $C$ 的一个二重覆盖上的 Heisenberg 共形块之间给出一个通道。该构造建立在 Zamolodchikov 的分支自由场构造 [373] 之上，不过此处所涉及的归一化关联函数如今处处正则，因此无需再考虑 Virasoro 原初场的额外插入。}

\section{Virasoro 与 Heisenberg 共形块}

抽象地接近共形块的一种方式，是把它们视为满足手征 Ward 恒等式的关联函数组（例如见 [139]、[342]）。这一进路允许把共形块空间 $\Conf(C,\mathcal{V})$ 定义为一个典范定义的向量空间，其中不涉及裤子分解之类的任意选择：它只依赖于一个顶点代数 $\mathcal{V}$ 与一个黎曼曲面 $C$ 的数据，并指定插入在穿孔处的原初场（如果有）。我们将只考虑 Virasoro 与 Heisenberg 顶点代数的情形。

对固定的常数 $c\in\C$，Virasoro 顶点代数 $\mathrm{Vir}_c$ 由单个场 $T$ 生成，它满足如下算子乘积展开（OPE）：
\begin{equation}
T(p)T(q)=\frac{c/2}{(z(p)-z(q))^4}+\frac{2T(q)}{(z(p)-z(q))^2}+\frac{\partial_{z(q)}T(q)}{z(p)-z(q)}+\text{regular terms}.
\tag{26.1}
\end{equation}
原初场 $W_h$（$h\in\C$）由如下条件定义：
\[
T(p)W_h(q)=\frac{hW_h(q)}{(z(p)-z(q))^2}+\frac{\mathrm{reg}}{z(p)-z(q)}.
\]
Heisenberg 顶点代数 Heis 由单个场 $J$ 生成，其 OPE 为：
% [原文 p.416]
\begin{equation}
J(p)J(q)=\frac{1}{(z(p)-z(q))^2}+\text{regular terms}.
\tag{26.2}
\end{equation}
原初场 $V_\alpha$（$\alpha\in\C$）由如下条件定义：
\begin{equation}
J(p)V_\alpha(q)=\frac{\alpha V_\alpha(q)}{z(p)-z(q)}+\text{regular terms},
\tag{26.3}
\end{equation}
并且可以验证 Heisenberg 顶点代数包含一个中心荷 $c=1$ 的 Virasoro 顶点代数（见 [139, Ex. 2.5.9]）。

现在设 $C$ 是由复图册 $(U_i,z_i)$ 覆盖的黎曼曲面，并给定点 $p_1,\dots,p_n\in C$（$n\in\N$ 固定，且 $p_i\in U_i$）。Heisenberg 共形块 $\Psi\in\Conf(C,\mathrm{Heis})$ 定义为对每个 $n$ 给出的关联函数组：
\[
\langle J(p_1)^{z_1}\cdots J(p_n)^{z_n}\rangle_\Psi: U_1\times U_2\times\cdots\times U_n\to\C,
\]
它们在除对角线 $p_i=p_j$ 以外全纯，而对角线处的奇性由 (26.2) 控制。现在取顶点代数 $\mathrm{Vir}_{c=1}\otimes\mathrm{Heis}$，并类似地定义 $c=1$ 的 Virasoro 共形块 $\Psi\in\Conf(C,\mathrm{Vir}_{c=1}\otimes\mathrm{Heis})$，作为关联函数组：
\[
\langle T(p_1)^{z_1}\cdots T(p_n)^{z_n}J(q_1)^{w_1}\cdots J(q_m)^{w_m}\rangle_\Psi: U_1\times U_2\times\cdots\times U_n\to\C,
\]
它们在点 $p_i=p_j$ 或 $q_i=q_j$ 之外全纯，满足 OPE (26.1) 与 (26.2)，以及某些坐标变换规则。

插入原初场的共形块空间 $\Conf(C,\mathrm{Heis};V_{\alpha_1}(q_1)\cdots V_{\alpha_k}(q_k))$ 定义为如下关联函数组所成的空间：
\[
\langle J(p_1)^{z_1}\cdots J(p_n)^{z_n}V_{\alpha_1}(q_1)\cdots V_{\alpha_k}(q_k)\rangle_\Psi: U_1\times U_2\times\cdots\times U_n\to\C,
\]
其中 $J$ 的插入具有与之前相同的 OPE 与坐标变换，但现在当任何 $p_i$ 遇上任何 $q_j$ 时，按 (26.3) 的要求会出现额外的一阶极点。

\section{非阿贝尔化映射}

Virasoro 块通常难以计算；不过可以利用特定的递推关系（见 [374]、[87]），或利用 AGT 对应 [11] 所提供的表示。计算它们的另一种可能途径是把问题化归为较易得到的 Heisenberg 共形块。这关系到所谓的自由场构造，它产生一类非常特殊的 Virasoro 块。具体地说，对常数 $Q\in\C$，把 $T=\frac12J^2+Q\partial J$ 的表达式代入 Heisenberg 关联函数，便给出一个映射
% [原文 p.417]
\[
\Conf(C,\mathrm{Heis})\to\Conf(C,\mathrm{Vir}_{c=1+6Q^2})
\]
把共形块映到共形块。当 $c=1$ 时，这一技巧的一个变体称为分支自由场构造，由 Zamolodchikov 在 [373] 中给出。对该构造，考虑一个分歧二重覆盖 $\pi:\tilde C\to C$，以及 $\tilde C$ 上的 Heisenberg 场 $\tilde J$。令
\[
\tilde J^{(-)}:=\frac{1}{\sqrt2}\left(\tilde J^{(1)}-\tilde J^{(2)}\right)
\]
表示 $\tilde C$ 的两叶上插入的反不变组合；它在差一个 $\Z_2$ 作用 $\tilde J^{(-)}\to-\tilde J^{(-)}$ 的意义下给出 $C$ 上良定义的算子。于是，把 $T:=\frac12(\tilde J^{(-)})^2$ 代入 Heisenberg 关联函数，便给出 $C$ 上的 Virasoro 关联函数，而诱导出的 Virasoro 块在覆盖的分支点 $b_1,\dots,b_k$ 处有额外奇性。这些奇性可视为权为 $h=\frac{1}{16}$ 的 Virasoro 原初场 $W_h(b_i)$ 的插入，因此存在一个良定义的映射
\[
\Conf(\tilde C,\mathrm{Heis})\to\Conf(C,\mathrm{Vir}_{c=1};W_{\frac{1}{16}}(b_1)\cdots W_{\frac{1}{16}}(b_k)).
\]

谱网络可用来修改上面概述的分支自由场构造，使得分支点处的 $W_{\frac{1}{16}}$ 插入不出现，而且别处也不产生额外奇性。这确实可以通过在 $\tilde C$ 上的关联函数中插入一个额外算子来实现。我们下面概述 Hao 与 Neitzke [188] 的这一构造。

对光滑二重覆盖 $\pi:\tilde C\to C$，固定一个从属于 $\pi$ 的 $\mathfrak{gl}_2$-谱网络 $\mathcal{W}$ 的一般选择（它在覆盖 $\pi$ 的每个分支点上终止三条弧），并固定 $C$ 上的一个自旋结构。然后考虑谱网络算子
\[
E(\mathcal{W})=\exp\left[\frac{1}{2\pi\ii}\int_{\mathcal{W}}\psi_+(z^{(+)})\psi_-(z^{(-)})\mathrm{d}z\right],
\]
其中 $\psi_\pm$ 是由 Heisenberg 场 $\tilde J$ 经费米化构造出的自由费米子，而 $z^{(+)},z^{(-)}\in\tilde C$ 是点 $z\in C$ 在二重覆盖 $\tilde C$ 上的两个原像。

[188] 中算出：当把 $E(\mathcal{W})$ 插入 $\tilde C$ 上的关联函数后，带有 $T(z)$ 的归一化关联函数处处正则，即使在分支点处也如此。因此，插入 $E(\mathcal{W})$ 去掉了分支点处不需要的 $W_{\frac{1}{16}}$ 插入，且不产生任何进一步奇性。此外，由于插入 $E(\mathcal{W})$ 后的 0 点函数发散，算子 $E(\mathcal{W})$ 被一个重整化版本 $E_{\mathrm{ren}}(\mathcal{W})$ 替代，以使 0 点函数有限。这一重整化引入对每个分支点附近局部坐标的反常依赖，其权为 $-\frac{1}{16}$，从而抵消这些点处权为 $\frac{1}{16}$ 的插入。

总之，共形块空间之间存在一个良定义的线性映射
\begin{equation}
\mathcal{F}_{\mathcal{W}}:\Conf(\tilde C,\mathrm{Heis})\to\Conf(C,\mathrm{Vir}_{c=1}\otimes\mathrm{Heis}),
\tag{26.4}
\end{equation}
% [原文 p.418]
称为共形块的非阿贝尔化映射；并且，对 $\tilde C$ 上的 Heisenberg 块 $\tilde\Psi$，相应块 $\mathcal{F}_{\mathcal{W}}(\tilde\Psi)$ 中 Virasoro 生成元的 1 点函数由
\[
\langle T(z)\rangle_{\mathcal{F}_{\mathcal{W}}(\tilde\Psi)}=\frac14\langle\tilde J(z^{(1)})-\tilde J(z^{(2)})^2E_{\mathrm{ren}}(\mathcal{W})\rangle_{\tilde\Psi}
\]
给出。

共形块空间 $\Conf(\tilde C,\mathrm{Heis})$ 上有一个由圈算子 $L_\gamma$ 给出的作用，这些算子以 $\tilde C$ 上的环路 $\gamma$ 为标记。这些圈算子生成交换的 skein 代数 $\mathrm{Sk}_{-1}(\tilde C,\GL(1))$，而圈算子代数对偶地是（扭）平坦联络的模空间 $\mathcal{M}(\tilde C,\GL(1))$ 上的函数代数。现在，对 $\mathcal{M}(\tilde C,\GL(1))$ 中的点 $X$，取它关于 $\tilde C$ 上一组固定的 $A$、$B$ 循环基的分解为 $X=(\ee^{x},\ee^{y})$。一般而言，我们可以构造一族由 $\tilde g$ 个连续参数参数化的线性无关的 Heisenberg 块
\[
\tilde\Psi_a\in\Conf(\tilde C,\mathrm{Heis}),
\]
其中 $a=(a_1,\dots,a_{\tilde g})\in\C^{\tilde g}$（见 [188, Sect. 6]）。特别地，我们把本征块 $\tilde\Psi_{x,y}\in\Conf(\tilde C,\mathrm{Heis})$ 写成 $\tilde\Psi_a$ 的线性组合：
\[
\tilde\Psi_{x,y}=\sum_{n\in\Z^{\tilde g}}\exp\left(-\frac{(x+2\pi\ii n)\cdot y}{2\pi\ii}\right)\tilde\Psi_{a=x+2\pi\ii n}.
\]
为了构造本征块 $\Psi\in\Conf(C,\mathrm{Vir}_{c=1}\otimes\mathrm{Heis})$，人们利用映射 $\mathcal{F}_{\mathcal{W}}$（来自 (26.4)）关于圈算子的协变性。这样我们得到一族本征块
\[
\Psi^{\mathcal{W}}_{x,y}:=\mathcal{F}_{\mathcal{W}}(\tilde\Psi_{x,y}).
\]

把共形块的抽象语言翻译到自由费米子场论：本征块 $\tilde\Psi_{x,y}$ 中 $\tilde C$ 上的关联函数，应理解为与 $\tilde C$ 上由具有和乐 $(\ee^{x},\ee^{y})$ 的复背景规范场扭的手征自由费米子理论相关。另一方面，本征块 $\Psi^{\mathcal{W}}_{x,y}$ 中 $C$ 上的关联函数，同样应当与 $C$ 上由具有和乐 $\lambda=\mathcal{F}^{\flat}_{\mathcal{W}}((\ee^{x},\ee^{y}))$ 的复背景规范场扭的两个手征自由费米子理论相关。从此观点看，(26.4) 中的非阿贝尔化映射 $\mathcal{F}_{\mathcal{W}}$ 勾勒出这两个场论之间的一个通道，也就是说，两个费米子行列式并不严格相等，但在把算子 $E_{\mathrm{ren}}(\mathcal{W})$ 插入到 $C$ 上的理论中之后，它们确实会相差一个常数而相等。

% [Tutor Note: 1 点函数公式照原文作 $\frac14\langle\tilde J(z^{(1)})-\tilde J(z^{(2)})^2E_{\mathrm{ren}}(\mathcal{W})\rangle$，但平方号只落在第二项上，与上下文不符：由 $T=\frac12(\tilde J^{(-)})^2$ 与 $\tilde J^{(-)}=\frac{1}{\sqrt2}(\tilde J(z^{(1)})-\tilde J(z^{(2)}))$ 立即得到 $T=\frac14(\tilde J(z^{(1)})-\tilde J(z^{(2)}))^2$，故平方应作用于整个差，正确形式为 $\frac14\langle(\tilde J(z^{(1)})-\tilde J(z^{(2)}))^2E_{\mathrm{ren}}(\mathcal{W})\rangle$。原文漏了外层大括号（照译并在此指出）；核对方法：展开 $T=\frac12(\tilde J^{(-)})^2$ 即得。]

% [Tutor Note: 本块另有一处归一化上的疑点（不改正文，供 Lead 与作者核对）：按 (26.1) 中“四阶极点系数 $=c/2$”与 (26.2) 中系数为 $1$ 的 OPE，对 $T=\frac12J^2+Q\partial J$ 直接计算 $\langle T(z)T(w)\rangle$ 得 $c/2=\frac12+6Q^2$，即 $c=1+12Q^2$，而原文所给映射是 $\mathrm{Vir}_{c=1+6Q^2}$。相差的因子 2 通常来自对 $J$ 的归一化（或 Liouville 文献中 $c=1+6Q^2$、$Q=b+b^{-1}$ 的惯例）。核对时以 (26.1)、(26.2) 的 OPE 系数为准；此点已在回报中列为存疑。]
